%% JSTARS submission (journal version of the IGARSS manuscript)
%% Compile: pdflatex main_jstars && pdflatex main_jstars   (IEEEtran, journal option)
%%
%% Every number in this manuscript is reproduced by the commands listed in
%% 13_结果清单_已核验.md, which recomputes them from exp/runs.csv. Figures come
%% from scripts/make_jstars_figures.py and read the same log.
\documentclass[journal]{IEEEtran}

\usepackage{amsmath}
\usepackage{graphicx}
\usepackage{booktabs}
\usepackage{url}
\usepackage[hidelinks]{hyperref}

% Allow hyphenation of the compound terms this paper uses repeatedly.
% \hyphenation accepts lowercase letters and hyphens only -- digits are
% rejected with "Not a letter", so dataset names containing numerals (e.g.
% TimeSen2Crop) cannot appear here.
\hyphenation{mask-aware sep-a-ra-ble breizh-crops time-sen-crop}

\begin{document}

\title{Masking Is Not One Mechanism: Separating the Attention Operator from
Its Input Channels in Crop Classification}

\author{Author Name%
\thanks{Manuscript received ...; revised ... .}%
\thanks{The author is with Department, Institution, City, Country (e-mail:
email@example.com).}}

\maketitle

\begin{abstract}
Masking unobserved steps in a temporal attention model is usually implemented as
one change, but it bundles two: the attention operator that suppresses missing
timesteps, and the input channels it arrives with---a mask indicator, elapsed
time and a calendar encoding. Because these vary together, the conventional
comparison against an unmasked model cannot say which one produced an observed
effect. We separate them with a capacity-matched control---identical
architecture and identical parameter count, differing only in whether the
operator is applied---and sweep observation removal from 0\% to 90\% on three
crop classification benchmarks (TimeSen2Crop, 15 classes; BreizhCrops, 7;
PASTIS, 18) and two attention families (query-as-parameter L-TAE and a standard
Transformer encoder). Three results follow. The operator's effect is not a
property of masking but of the benchmark: it is strongly negative on
TimeSen2Crop ($-0.0648$ at 90\% removal, $t=-20.0$), indistinguishable from zero
on BreizhCrops ($-0.0014$), and strongly positive on PASTIS ($+0.0569$,
$t=+13.0$). The input block is itself three mechanisms, not one: elapsed time is
harmful on two of the three benchmarks ($-0.0345$ and $-0.0122$ at zero
removal), the mask indicator is inert throughout, and the calendar drives the
BreizhCrops effect on its own. Shuffling each channel's values across timesteps
within a shard---same distribution, same per-shard constancy, alignment
destroyed---shows the calendar is read for its temporal content on both
benchmarks where it matters, by $+0.0621$ and $+0.0807$ at 90\% removal. On
BreizhCrops the decomposition is unchanged under the second attention family,
so the benchmark's effect belongs to the channels rather than to the operator
the comparison is named for.
\end{abstract}

\begin{IEEEkeywords}
Sentinel-2, crop classification, temporal attention, missing data, cloud gaps,
robustness, reproducibility
\end{IEEEkeywords}

\section{Introduction}

Optical satellite time series for agriculture are routinely incomplete. In the
Sentinel-2 archives used here, the clear-sky fraction varies by a factor of 2.2
across tiles (0.383 to 0.832), and it is strongly seasonal: pooled over the
TimeSen2Crop tiles it runs from 0.261 in February to 0.863 in August, and in a
single tile the December value can fall to 0.032. Any operational crop mapping
system must therefore decide what to do with observations it does not have.

Temporal attention models answer this in a now-standard way: mask the
unobserved steps before the softmax, so they cannot attract attention. The
rationale is straightforward, as an unavailable observation carries no
information and should not influence the output. Self-supervised approaches
extend the same idea, treating cloud occlusion as a natural mask~\cite{mae}.
The mechanism appears in widely used architectures~\cite{ltae,tsvit,alise}, and
recent work continues to build on it for irregular, unaligned
sampling~\cite{alise,shaping}.

Whether it helps has not been bounded: a review of 29 public benchmarks finds no
dedicated protocol for evaluating robustness to missing observations, so results
are measured under conditions that differ between papers~\cite{sitsreview}.

We ask a narrower and answerable question: \emph{under controlled degradation,
when does mask-aware attention help, and when does it hurt?} Answering it
requires separating two things that are normally changed together, and the
separation is the contribution:
\begin{enumerate}
\item A controlled degradation protocol with two orthogonal axes, per-sample
matched real-cloud controls, and hard guards against geographic leakage
(Sec.~\ref{sec:protocol}).
\item A capacity-matched decomposition separating the attention masking
operator from the input channels it is bundled with, applied across three
benchmarks and two attention families. The two components have opposite effects,
and the benchmarks disagree about which dominates, and in which direction
(Sec.~\ref{sec:decomp}).
\item A three-way split of the input block into mask indicator, elapsed time and
calendar encoding, each with a time-alignment shuffle control that distinguishes
a channel carrying temporal information from a channel merely present. The
block is not one mechanism: its parts have different signs, and the largest
effect on one benchmark is the calendar alone (Sec.~\ref{sec:channels}).
\item A within-sample sequence-length intervention showing the operator effect
has no trend with the number of surviving observations, and that the pattern
suggested by an observational stratification does not survive it
(Sec.~\ref{sec:trl}).
\end{enumerate}

Throughout, we report negative results as results. Two candidate third
benchmarks were acquired, measured and discarded (Sec.~\ref{sec:datasets}), and
one moderator we had previously reported does not survive a stricter
test. Both are reported rather than dropped.

\section{Related Work}

\subsection{Handling missing observations}
Work on incomplete satellite image time series (SITS) falls into three
families. Reconstruction-then-classify pipelines impute gaps before prediction;
a recent joint framework argues instead that full reconstruction is often
unnecessary and propagates error~\cite{jointfrp}. End-to-end approaches
tolerate gaps directly, either through masking~\cite{mae} or through encoders
designed for irregular, unaligned sampling~\cite{alise}. A third family
sidesteps optical gaps using SAR~\cite{sitsreview}. Our work belongs to the
second family and tests the assumption that underwrites it.

\subsection{Lightweight temporal attention}
The Lightweight Temporal Attention Encoder (L-TAE) introduced channel grouping
and query-as-parameter attention, reducing parameter counts by orders of
magnitude relative to the Temporal Attention Encoder~\cite{ltae}. U-TAE embeds
the same encoder in a U-Net for panoptic
segmentation~\cite{utae}. TSViT adds acquisition-time-specific positional
encodings within a factorized temporo-spatial transformer~\cite{tsvit}. Our
primary model is L-TAE augmented with masking, elapsed-time and calendar
encodings, a 256-parameter difference, which keeps the comparison attributable
to the mechanism rather than to capacity. For the same reason we add a standard
Transformer encoder~\cite{vaswani} as a second, structurally different attention
family rather than a larger model of the same kind.

\subsection{Candidate moderators, and padding}
Temporal resolution matters: accuracy on PASTIS-HD plateaus logarithmically as
input length grows from 1 to 35 steps~\cite{shaping}. Class-confusion studies
find that certain crop pairs are persistently confused irrespective of
architecture~\cite{eurosat}. We tested both as moderators of masking; neither
survives (Sec.~\ref{sec:trl}). Padding variable-length sequences to a common
length also measurably changes performance~\cite{padding}---a trap we fell into
ourselves (Sec.~\ref{sec:protocol}). Reference~\cite{mae} is a conference
abstract rather than a peer-reviewed paper, and is cited as such.

\section{Controlled Degradation Protocol}
\label{sec:protocol}

\subsection{Two orthogonal axes}
We distinguish two ways a time series becomes incomplete, which the literature
often conflates. \emph{Missingness} means the timestamp is retained but the
observation value is unavailable. \emph{Irregularity} means the timestamps
themselves are unevenly spaced. Only the first is manipulated here.
Irregularity is an intrinsic property of all three datasets (median
inter-acquisition interval 5--10 days; maximum 95 days in TimeSen2Crop).

\subsection{Degradation modes}
\emph{Random removal} removes a uniformly random subset, matching rate to the
original sequence length. \emph{Real cloud} drops observations flagged as cloud
or shadow by the dataset's own quality band. \emph{Matched random} is the
control for real cloud: it removes the \emph{same number} of observations from
each individual sample, uniformly at random. Matching is per sample rather than
per tile or globally, because cloudiness varies pixel to pixel; a tile-average
match would leave the arms differing in both structure and per-pixel
difficulty, and a global match would confound structure with which tiles happen
to be cloudy. We verify that the two arms have identical per-sample observation
counts (max absolute difference: 0).

\subsection{Time-channel alignment control}
At zero removal the elapsed-time and calendar channels are \emph{per-shard
constant vectors}: every sample of a shard carries the same values. An effect
measured for such a channel therefore admits two explanations. Under the
\emph{temporal} explanation the model reads the acquisition schedule or the
season from the values and their position in the sequence, so destroying the
alignment must destroy the effect. Under the \emph{presence} explanation the
effect comes from the vector existing at all---an extra learned per-timestep
bias, or a tag identifying the shard and with it the regional crop mix---and
alignment is irrelevant.

We separate these by permuting the named channel's values across timesteps
\emph{within each shard}. The value distribution is unchanged, the vector
remains per-shard constant, and only the alignment is broken. A channel whose
measured step collapses toward zero under permutation carries temporal
information; a channel whose step survives does not, and must not be described
as temporal. The control arms are trained exactly like the arms they control
for, and are reported under their own tag so that they never pool with them.

\subsection{Determinism}
Every stochastic function takes an explicit generator and is bit-reproducible
given a seed. This is enforced by tests rather than by convention, because the
benchmark's central claim is that its degradations are reproducible.

\subsection{Five traps that silently corrupt results}
Five failure modes biased results without raising an error.
\emph{Day-of-year} resets at 1 January, so differencing it produced a
$-357$-day interval. \emph{Padding} per-tile series to a split maximum gave
0--28.9\% padding in training against 0--3.3\% at test, manufacturing an
apparent 6.6-point masking advantage. \emph{Best-epoch selection} tripled
run-to-run variance ($\sigma\,0.040 \to 0.009$), and \emph{BatchNorm
recalibration} after weight averaging sampled 2 of 12 tiles, costing 11 points.
A \emph{bare \texttt{YYYYMMDD} date} parses as a float, silently discarding
68\% of a dataset.

A sixth trap belongs to the third benchmark's archive rather than to our code,
and cost more time than the other five together. The published description of
the PASTIS archive disagrees with the archive in three places: the band members
are named \texttt{S2\_<id>.npy} and \texttt{TARGET\_<id>.npy} rather than
\texttt{<id>.npy}; the reflectance is signed \texttt{int16} at $10^{4}$ scale
rather than \texttt{uint16}; and the per-patch acquisition dates are stored as a
dictionary keyed by time index rather than as a list. The third is the dangerous
one: reading the keys instead of the values yields plausible-looking integers,
and the dictionary's length still equals the sequence length, so a consistency
check on length passes. We now validate archive structure against the archive's
own member names before trusting any description of it.

\section{Experimental Setup}

\subsection{Datasets}
\label{sec:datasets}
The three benchmarks differ deliberately in scale, unit and sampling.
TimeSen2Crop~\cite{timesen2crop} provides 1{,}212{,}224 labelled pixels over 15
classes in 9 bands with 23--38 acquisitions, split into 12 training tiles,
one validation tile (33TWN) and two test tiles (33UVP, 33UXP).
BreizhCrops~\cite{breizhcrops} provides 601{,}319 parcels over 7 classes in 10
bands with 16--175 acquisitions, split by region into FRH01+02 (train), FRH03
(validation) and FRH04 (test). PASTIS~\cite{utae} provides 2{,}433 patches over
18 crop classes in 10 bands with 38--61 acquisitions, on the archive's own five
folds---folds 1--3 train, 4 validate, 5 test---which the archive constructs so
that neighbouring patches never fall in different folds, a stricter spatial
separation than either of the other two benchmarks has. Table~\ref{tab:data}
summarizes.

\begin{table}[t]
\centering
\caption{The three benchmarks. Removal rates are swept identically on all
three; the split is fixed by each dataset's own published partition and is
never chosen after seeing a score.}
\label{tab:data}
\footnotesize
\begin{tabular}{llll}
\toprule
 & TimeSen2Crop & BreizhCrops & PASTIS \\
\midrule
Unit            & pixel   & parcel  & pixel \\
Classes         & 15      & 7       & 18 \\
Bands           & 9       & 10      & 10 \\
Acquisitions    & 23--38  & 16--175 & 38--61 \\
Train samples   & 646{,}083 & 316{,}023 & 582{,}000 \\
Eval samples    & 253{,}468 & 120{,}140 & 198{,}400 \\
Split           & tiles   & regions & archive folds \\
Country         & Italy   & France  & France \\
\bottomrule
\end{tabular}
\end{table}

BreizhCrops parcels are grouped by acquisition-date signature so that every
evaluation shard has uniform length; without this, batching would reintroduce
the padding trap above. Our reconstructed parcel count (601{,}319) agrees with
the published 608{,}263.

Because our TimeSen2Crop tiles share one acquisition schedule, a within-tile
split would give the test set the same inter-observation intervals as training.
We hold out entire tiles, and a hard guard rejects any split placing the same
ground on both sides; this excludes the 2019 acquisition of a 2018 test tile.

\paragraph{Two datasets we acquired, measured and discarded.}
A reviewer of the conference version asked for a third benchmark under a
different geographic or agricultural regime, and that request shaped our
selection. Two non-European candidates were obtained and rejected on
measurement, not on availability.

A South African product carried a per-scene radiometric state that dominates the
crop signal. In a 38-dimensional NDVI representation the same crop differs
$1.58\times$ more between two spatial chunks ($1.087$ in $L_2$) than two
different crops do inside one chunk ($0.687$); a classifier separates crops
inside a chunk at $0.87$ and across chunks at $0.095$. Under that representation
cross-region transfer is not merely hard but unavailable at the stated feature
level, and the smoke-test macro-F1 of $0.16$--$0.23$ is consistent with it
against a nine-class chance level of $0.11$.

CONUS HLS+CDL carries genuinely consistent surface reflectance---atmospheric
correction, BRDF normalisation and Landsat/Sentinel-2 harmonisation---and
passed a structural gate built for exactly that property. It nevertheless did
not transfer across tiles: a random forest separates crops inside one tile at
$0.70$ and across tiles at $0.23$, and a 14k-parameter neural model reaches
macro-F1 $0.10$--$0.13$ after 150 epochs. Region-wise the failure is uneven---in
one irrigated cluster the cross-region and within-region scores coincide at
$0.63$---but selecting that cluster after seeing its score would be
cherry-picking, so the dataset was dropped whole.

The generalizable observation, which we offer as a methodological result in its
own right, is that \emph{the spatial splits reported in this literature are
within-region splits}. BreizhCrops holds out regions roughly 100~km apart inside
Brittany; TimeSen2Crop holds out neighbouring tiles inside the Po Valley;
PASTIS's folds are separated by at least 1~km within metropolitan France. None
of the openly downloadable benchmarks we examined supports a
continental-scale spatial holdout, and a paper that claims otherwise should
expect the claim not to replicate. PASTIS is French, exactly like BreizhCrops,
and therefore answers the request for a different agricultural regime only in
part. We state this as a limitation rather than obscure it: what the third
benchmark buys is not geography but a third independent configuration---a
different taxonomy, a different spatial unit and a different split
construction---in which to test whether the decomposition generalizes.

\subsection{Models}
MT-TAE is L-TAE~\cite{ltae} with three added mechanisms: mask-aware attention
(logits of unobserved steps set to $-\infty$ before the softmax), an
elapsed-time ($\Delta t$) embedding, and a calendar encoding.
\emph{Operator} and \emph{channels} are separate switches, and the whole
decomposition rests on that separation: the operator costs zero parameters,
while the three input channels cost 256 (14{,}672 versus 14{,}416 for plain
L-TAE and 14{,}672 for the capacity-matched control, which differs from MT-TAE
only in whether the operator is applied).

The second family is a standard Transformer encoder~\cite{vaswani} over the same
inputs, with the same two switches and the same channel definitions. Adding a
family is what tests whether a decomposition is a property of the masking
mechanism or of query-as-parameter attention specifically. We also report
TempCNN (64k), GRU (154k) and LSTM (204k) as capacity controls; these have no
masking operator and so cannot enter the decomposition, but they anchor the
clean-data ranking.

\subsection{Training}
30 epochs, batch size 512, AdamW with cosine decay, mixed precision. We average
the last ten epochs and recalibrate BatchNorm statistics. All results are mean
$\pm$ standard deviation over three seeds. Contrasts are paired by seed, which
is the design-correct test here: the same seed trains both arms of a contrast,
so the pairing removes the seed-to-seed variation that would otherwise enter the
standard error. Deltas are identical under either convention; only the test
statistic changes, and we report the unpaired value beneath each table for
reference.

\paragraph{A caveat on absolute accuracy.}
TimeSen2Crop's authors capped sampling at 15{,}000 parcels per tile--class,
which makes its test set artificially balanced. Absolute macro-F1 values there
should not be compared with figures reported under the true crop distribution.
Our comparisons are relative and measured within a fixed evaluation set.

\section{Results}
\label{sec:results}

\subsection{Clean-data baselines}
Table~\ref{tab:baselines} gives clean-data performance. The ranking broadly
tracks parameter count---the three larger models hold the top three places and
the two 14k-parameter attention models the bottom two---which serves as a check
that the protocol does not distort capacity comparisons. It is not strictly
monotone: GRU and LSTM are 0.0009 apart with standard deviations of 0.012 and
0.007, so we do not claim an ordering between them. On the two benchmarks where
the attention family is available, clean accuracy is comparable across datasets
despite their very different units and label sets (L-TAE 0.8365 on
TimeSen2Crop, 0.7129 on BreizhCrops, 0.6733 on PASTIS), so no benchmark is
unreachable for a model of this size.

A secondary result deserves note. The two smallest models are by a wide margin
the best calibrated: MT-TAE reaches an expected calibration error of 0.0180 and
L-TAE 0.0187, against 0.0378--0.0410 for the three larger models. MT-TAE is
14$\times$ smaller than the LSTM, yet its calibration error is under half as
large.

\begin{table}[t]
\centering
\caption{Clean-data baselines on TimeSen2Crop (15 classes, 253{,}468 evaluation
samples; the two test tiles hold 263{,}887 pixels, of which 253{,}468 survive
the split guards and enter the metric). Mean over three seeds.}
\label{tab:baselines}
\begin{tabular}{lrrrr}
\toprule
Model & macro-F1 & s.d. & ECE & Params \\
\midrule
GRU     & 0.8987 & 0.0115 & 0.0388 & 154{,}385 \\
LSTM    & 0.8978 & 0.0066 & 0.0410 & 204{,}049 \\
TempCNN & 0.8555 & 0.0066 & 0.0378 & 64{,}145 \\
MT-TAE  & 0.8406 & 0.0059 & \textbf{0.0180} & 14{,}672 \\
L-TAE   & 0.8365 & 0.0042 & \textbf{0.0187} & 14{,}416 \\
\bottomrule
\end{tabular}
\end{table}

\subsection{The operator and its input channels have separable effects}
\label{sec:decomp}

The masking mechanism bundles two changes that standard comparisons vary
together: the attention masking operator itself (zero parameters) and the input
channels it arrives with---a mask indicator, elapsed time and a calendar
encoding (256 parameters). Comparing L-TAE against a masked variant therefore
changes capacity and operator simultaneously, and cannot attribute the result to
either. We add a capacity-matched control, \texttt{mttae\_nomaskop}, which has
identical architecture and \emph{identical} parameter count (14{,}672) to
MT-TAE and differs only in whether the masking operator is applied.

\begin{table}[t]
\centering
\caption{Capacity-matched decomposition, $\Delta$ macro-F1 with paired $t$
statistics in parentheses. Operator effect compares \texttt{mttae} against
\texttt{mttae\_nomaskop} (same 14{,}672 parameters); input-channel effect
compares \texttt{mttae\_nomaskop} against L-TAE. The operator's sign is
dataset-specific; the channels' sign is not, but its size is.}
\label{tab:decomp}
\footnotesize
\begin{tabular}{lcccccc}
\toprule
 & \multicolumn{2}{c}{TimeSen2Crop} & \multicolumn{2}{c}{BreizhCrops}
 & \multicolumn{2}{c}{PASTIS} \\
\cmidrule(lr){2-3}\cmidrule(lr){4-5}\cmidrule(lr){6-7}
Removed & Op. & Chan. & Op. & Chan. & Op. & Chan. \\
\midrule
0\%  & $+0.0000$ (0.0)   & $+0.0041$ (0.8)   & $+0.0000$ (0.0)   & $-0.0058$ (-3.9)  & $+0.0000$ (0.0)   & $-0.0070$ (-5.3) \\
10\% & $-0.0006$ (-0.2)  & $+0.0105$ (1.7)   & $+0.0013$ (0.9)   & $-0.0085$ (-3.9)  & $+0.0077$ (2.3)   & $-0.0123$ (-4.6) \\
20\% & $-0.0055$ (-1.0)  & $-0.0002$ (-0.0)  & $-0.0017$ (-0.5)  & $-0.0105$ (-9.7)  & $+0.0137$ (4.9)   & $-0.0188$ (-4.3) \\
30\% & $-0.0163$ (-6.2)  & $+0.0130$ (3.3)   & $-0.0011$ (-0.4)  & $-0.0178$ (-4.4)  & $+0.0117$ (2.4)   & $-0.0146$ (-2.3) \\
50\% & $-0.0178$ (-3.8)  & $+0.0061$ (0.7)   & $-0.0004$ (-0.2)  & $-0.0165$ (-3.7)  & $+0.0134$ (3.0)   & $-0.0120$ (-2.7) \\
70\% & $-0.0264$ (-3.9)  & $-0.0038$ (-0.5)  & $+0.0049$ (1.5)   & $-0.0107$ (-2.7)  & $+0.0259$ (14.3)  & $-0.0200$ (-11.0) \\
90\% & $-0.0648$ (-20.0) & $+0.0031$ (3.0)   & $-0.0014$ (-0.6)  & $+0.0178$ (8.8)   & $+0.0569$ (13.0)  & $-0.0692$ (-8.9) \\
\bottomrule
\end{tabular}
\end{table}

\begin{figure*}[t]
\centering
\includegraphics[width=0.95\textwidth]{figs/jstars_fig1_decomposition.pdf}
\caption{Capacity-matched decomposition on the three benchmarks. Within a panel
both curves come from arms with identical architecture and identical parameter
count, so each isolates one component. The operator's curve changes sign
between panels while the channel curve does not: on TimeSen2Crop the operator
alone drives a monotone decline, on BreizhCrops it is flat and the channels
carry the pattern, and on PASTIS the two have swapped signs again.}
\label{fig:decomp}
\end{figure*}

The decomposition (Table~\ref{tab:decomp}, Fig.~\ref{fig:decomp}) is
unambiguous, and the three benchmarks disagree about the operator.

On TimeSen2Crop the operator alone is responsible for the entire effect. It is
neutral below 20\% removal, becomes significantly negative at 30\%
($-0.0163$, $t=-6.2$), and reaches $-0.0648$ ($t=-20.0$) at 90\%; the 256
parameters contribute nothing systematic at any rate. On BreizhCrops the
assignment inverts: the operator is inert at every rate (largest $|t| = 1.5$),
while the channels alone produce the entire pattern---significantly negative
through the mid-range ($-0.0178$ at 30\%, $t=-4.4$) and significantly positive
at 90\% ($+0.0178$, $t=+8.8$). On PASTIS the operator is the beneficial
component, significantly so from 10\% removal onward and reaching $+0.0569$
($t=+13.0$) at 90\%, while the channels are uniformly harmful and reach
$-0.0692$ ($t=-8.9$).

The three operator effects at 90\% removal are therefore $-0.0648$, $-0.0014$
and $+0.0569$. Inconsistency of this kind is normally a weakness to be
explained away. Here it is the argument: a comparison that changes both
components at once reports a single number whose sign is determined by which
component happens to dominate on the benchmark at hand, and which component
dominates is not a stable property of masking.

\subsection{The control is validated where it must be}
At zero removal the operator effect is exactly $0.0000$ on all three
benchmarks. This is the expected result and it is not
trivial: with no observation masked, masking the attention logits is an
identity transform, a fact we also verified by construction
(Sec.~\ref{sec:protocol}). The three zeros therefore confirm that the control
isolates the operator rather than some incidental difference in initialisation
or training. Any later divergence between the arms is attributable to the
operator acting on genuinely missing steps.

The equality is exact rather than approximate. At rate zero the two degradation
specifications compile to the same split, and the same seed trains the same
model, so the paired per-seed differences are $[0.0, 0.0, 0.0]$ to machine
precision on every benchmark. A control whose zero is exact is a stronger
statement than one whose zero is within noise.

\subsection{What the conventional comparison would have reported}
Comparing a masked model against an unmasked one at 90\% removal gives
$-0.0618$ on TimeSen2Crop, $+0.0164$ on BreizhCrops and $-0.0122$ on PASTIS.
The decomposition shows that on the first dataset $-0.0648$ of the difference is
the operator and $+0.0031$ is the channels, which nearly cancel; on the second,
$-0.0014$ is the operator and $+0.0178$ is the channels; on the third, $+0.0569$
and $-0.0692$ again nearly cancel. On two of the three benchmarks the reported
sign belongs to a component the comparison was not designed to test.

\subsection{Inside the channel block: three mechanisms, not one}
\label{sec:channels}

Calling the 256 parameters ``the input-channel effect'' still treats three
different things as one. The channel chain adds them one at a time with the
attention operator held off throughout, so that each step is measured against
the arm directly beneath it and no step carries the effect of its neighbours.

\begin{table}[t]
\centering
\caption{Channel decomposition, $\Delta$ macro-F1 with paired $t$ statistics.
Each channel is added on top of the ones above it, with the attention operator
off in every arm, so these are channel effects alone. The three channels have
different signs, and the two largest nearly cancel on BreizhCrops---which is
why the block read as a single weak effect.}
\label{tab:channels}
\scriptsize
\begin{tabular}{llrrr}
\toprule
Dataset & Removed & $+$mask & $+\Delta t$ & $+$calendar \\
\midrule
TimeSen2Crop & 0\%  & $+0.0029$ (1.1) & $+0.0055$ (0.5) & $-0.0044$ (-0.8) \\
             & 90\% & $+0.0006$ (0.1) & $+0.0003$ (0.1) & $+0.0021$ (0.8) \\
\addlinespace
BreizhCrops  & 0\%  & $-0.0050$ (-4.1) & $-0.0345$ (-32.7) & $+0.0337$ (83.6) \\
             & 90\% & $+0.0006$ (2.0) & $-0.0158$ (-16.3) & $+0.0330$ (16.8) \\
\addlinespace
PASTIS       & 0\%  & $+0.0015$ (0.5) & $-0.0122$ (-2.3) & $+0.0037$ (1.5) \\
             & 90\% & $+0.0013$ (0.3) & $-0.0166$ (-4.0) & $-0.0539$ (-6.6) \\
\bottomrule
\end{tabular}
\end{table}

\begin{figure*}[t]
\centering
\includegraphics[width=0.95\textwidth]{figs/jstars_fig2_channels.pdf}
\caption{The three channels of the input block, measured one at a time with the
attention operator off. Elapsed time costs accuracy on the two benchmarks where
anything happens at all and is inert on the third; the calendar is the largest
single effect on BreizhCrops and the largest on PASTIS, with opposite signs; the
mask indicator never produces a systematic effect. The non-monotone block-level
pattern reported for BreizhCrops is the sum of the two large, opposing curves.}
\label{fig:channels}
\end{figure*}

Three findings follow, and only the first is uniform across datasets
(Table~\ref{tab:channels}, Fig.~\ref{fig:channels}).

\emph{Elapsed time is harmful wherever it does anything.} On BreizhCrops it
costs $-0.0345$ macro-F1 at zero removal ($t=-32.7$, the largest single
statistic in this paper) and remains negative at every rate. On PASTIS it costs
$-0.0122$ at zero removal and grows to $-0.0166$ at 90\%. On TimeSen2Crop it is
inert. A negative effect at \emph{zero} removal is worth pausing on: nothing is
missing, so the channel is not telling the model about missingness. It is
supplying an inter-observation interval on clean data, and the model does worse
for it.

\emph{The calendar carries the BreizhCrops effect on its own.} At zero removal
it contributes $+0.0337$ ($t=+83.6$) and at 90\% $+0.0330$ ($t=+16.8$). Since
elapsed time contributes $-0.0345$, the block's net effect is small and
non-monotone, which is exactly the pattern the conference version reported and
could not explain.

\emph{The mask indicator is inert.} Across all three benchmarks and all rates
its step never exceeds $0.005$ in magnitude and its largest $|t|$ is 4.1. This
is the component the mechanism is named for.

On PASTIS the calendar's sign reverses, reaching $-0.0539$ ($t=-6.6$) at 90\%
removal. The next section separates the two readings of that reversal.

\subsection{Does the calendar encode time, or is it merely present?}
\label{sec:align}

At zero removal both time channels are per-shard constants, so their effects
admit the two explanations set out in Sec.~\ref{sec:protocol}. We permute each
channel's values across timesteps within each shard and re-measure the same
step. Table~\ref{tab:shuffle} reports the real step, the permuted step, and
their difference---the part of the effect that depends on alignment.

\begin{table}[t]
\centering
\caption{Alignment control. \emph{Real} is the channel's measured step;
\emph{shuffled} is the same step with that channel's values permuted across
timesteps within each shard. The difference is the part of the effect that
depends on alignment: large positive means the channel is read for its temporal
content. Deltas are paired by seed, $t$ in parentheses.}
\label{tab:shuffle}
\scriptsize
\begin{tabular}{llrrr}
\toprule
Dataset & Channel & Removed & Real & Shuffled \\
\midrule
BreizhCrops & $+\Delta t$ & 0\%  & $-0.0345$ (-32.7) & $-0.0288$ (-13.9) \\
            &             & 90\% & $-0.0158$ (-16.3) & $-0.0182$ (-5.0) \\
            & $+$calendar & 0\%  & $+0.0337$ (83.6)  & $-0.0284$ (-6.1) \\
            &             & 90\% & $+0.0330$ (16.8)  & $-0.0210$ (-3.4) \\
\addlinespace
PASTIS      & $+\Delta t$ & 0\%  & $-0.0122$ (-2.3)  & $-0.0257$ (-9.1) \\
            &             & 90\% & $-0.0166$ (-4.0)  & $-0.0342$ (-7.2) \\
            & $+$calendar & 0\%  & $+0.0037$ (1.5)   & $-0.0357$ (-6.9) \\
            &             & 90\% & $-0.0539$ (-6.6)  & $-0.1346$ (-10.7) \\
\bottomrule
\end{tabular}
\end{table}

\begin{figure*}[t]
\centering
\includegraphics[width=0.62\textwidth]{figs/jstars_fig3_alignment.pdf}
\caption{The calendar channel as measured and with its values permuted across
timesteps within each shard. The value distribution and the per-shard
constancy are identical in the two arms; only the alignment differs. On both
benchmarks the permuted channel is substantially worse, so the calendar is read
for its temporal content. On PASTIS the measured channel is still below the
no-channel baseline, which places the benchmark's disagreement in how the
information is used rather than in whether it exists.}
\label{fig:align}
\end{figure*}

The calendar's alignment matters on both benchmarks (Table~\ref{tab:shuffle},
Fig.~\ref{fig:align}). Permuting it costs $+0.0621$ macro-F1 relative to the
measured channel at zero removal on BreizhCrops and $+0.0539$ at 90\%; on PASTIS
the corresponding gaps are $+0.0394$ and $+0.0807$. Destroying the alignment
does not merely remove the channel's benefit, it makes the channel worse than
not having it at all ($-0.0284$ and $-0.0357$ at zero removal), which is what
one expects if the model is trying to read a schedule that no longer exists.

This settles the PASTIS question and sharpens it. The calendar there is
\emph{not} a mere shard fingerprint: a fingerprint would survive permutation,
and permutation instead costs $0.0807$ at 90\%. The channel carries real
temporal information on both benchmarks. What differs is what the model does
with it---on BreizhCrops the information is worth more than the channel costs,
on PASTIS it is worth less.

Elapsed time behaves differently, and inconsistently between the two
benchmarks. On BreizhCrops its harm is alignment-independent: the permuted step
is the same as the measured one ($-0.0288$ against $-0.0345$ at zero removal),
so the cost is not that the model is misled by a schedule it cannot use. It is
the cost of the vector being there at all. On PASTIS the permuted step is
\emph{worse} than the measured one at every rate, so there the alignment does
carry information, and the channel is harmful despite being informative. Both
readings rule out the account in which a time channel helps the model track
irregular sampling: on these benchmarks, adding elapsed time does not help even
when it is correctly aligned.

\subsection{A second attention family}
\label{sec:arch}

An objection to everything above is that the decomposition may be a property of
query-as-parameter attention rather than of masking. We therefore repeat the
BreizhCrops decomposition with a standard Transformer encoder~\cite{vaswani}
over the same inputs, the same two switches and the same channel definitions.

\begin{figure*}[t]
\centering
\includegraphics[width=0.62\textwidth]{figs/jstars_fig4_architecture.pdf}
\caption{The same decomposition under a second attention family, on
BreizhCrops. Under both families the operator is indistinguishable from zero at
every rate and the channels carry the pattern. The benchmark's effect belongs
to the input channels, not to the operator the conventional comparison is named
for.}
\label{fig:arch}
\end{figure*}

The result (Fig.~\ref{fig:arch}) is that the benchmark's attribution is
unchanged. The operator remains inert ($+0.0048$, $t=+2.5$, at 90\%; largest
$|t|$ over all rates 3.2), while the channels remain the component that
matters. The channel-level structure replicates as well: at zero removal the
calendar contributes $+0.0316$ and elapsed time $-0.0382$, against $+0.0337$
and $-0.0345$ for L-TAE. Two architectures with different attention
parameterisations, trained separately, put the same two numbers in the same two
places.

What this does \emph{not} establish is that the operator's cross-dataset
variation is architecture-independent. On BreizhCrops the operator is inert
under both families, which is consistent with its being inert there for reasons
that have nothing to do with architecture; but the two benchmarks where the
operator does something were only measured under L-TAE. Extending the second
family to TimeSen2Crop and PASTIS is the natural next test, and we flag it as
open rather than infer it.

\subsection{A controlled sequence-length intervention}
\label{sec:trl}

Our initial account attributed the cross-dataset disagreement to sequence
length, by stratifying test shards on their natural acquisition count. That
stratification is observational: long and short sequences are not randomly
assigned, and differ in crop composition, geography and seasonality at once.

We replace it with a within-sample manipulation. Taking the same parcels, we
subsample each to exactly $N \in \{20, 30, 45, 60, 80\}$ observations. The
per-sample permutation is drawn from a seed independent of $N$, so the subsets
are \emph{nested}: the $N=20$ set is contained in the $N=30$ set, and adding
observations is the only thing that changes.

\begin{table}[t]
\centering
\caption{Controlled sequence-length intervention. Same parcels, nested
observation subsets, capacity-matched arms.}
\label{tab:trl}
\footnotesize
\begin{tabular}{rrrr}
\toprule
$N$ kept & Operator off & Operator on & Difference \\
\midrule
20 & $0.6758 \pm 0.0060$ & $0.6797 \pm 0.0054$ & $+0.0038$ \\
30 & $0.6867 \pm 0.0034$ & $0.6894 \pm 0.0041$ & $+0.0027$ \\
45 & $0.6940 \pm 0.0016$ & $0.6986 \pm 0.0083$ & $+0.0046$ \\
60 & $0.7080 \pm 0.0074$ & $0.7074 \pm 0.0060$ & $-0.0006$ \\
80 & $0.7164 \pm 0.0034$ & $0.7170 \pm 0.0015$ & $+0.0006$ \\
\bottomrule
\end{tabular}
\end{table}

The operator effect does not trend with $N$ (Table~\ref{tab:trl}; slope
$-0.00006$ per added observation, all differences within noise), though absolute
accuracy rises steadily from 0.676 to 0.716. The manipulation works; the
sequence-length account does not survive it.

\subsection{Robustness to degradation structure}
Real cloud is not uniformly distributed: over the TimeSen2Crop tiles the
clear-sky fraction runs from 0.261 in February to 0.863 in August, and the
tile-level December spread alone is 0.032 to 0.988. We therefore repeat the
comparison against a control removing the
\emph{same number} of observations per sample, uniformly at random; the masking
deficit differs between arms by only $+0.0020$. The control also surfaces an
unanticipated effect---real cloud is harder than matched removal for the three
larger models ($-0.003$ to $-0.012$) but easier for the two attention models
($+0.026$, $+0.028$)---which we report untested.

\subsection{Mechanism}
\label{sec:mech}
Re-measured on the capacity-matched pair, the mechanism attributes exactly.
The operator alone raises within-class spread of the pooled representation
monotonically with removal: the ratio of masked to unmasked spread is exactly
$1.000$ at 0\% removal and rises through $1.210$ at 30\% to $\mathbf{1.631}$ at
90\%. Simultaneously it concentrates attention, effective attended observations
falling to $1.62$ at 90\% while its capacity-matched control holds at $2.73$.

The value of exactly $1.000$ at zero removal is the identity result measured
rather than argued: with nothing masked, two architectures differing only in
whether that operator is applied produce identical representations. On
TimeSen2Crop the chain is complete: masking concentrates attention onto fewer
survivors, the representation of a given class spreads, and accuracy falls.

That chain is stated for TimeSen2Crop deliberately. On BreizhCrops the operator
does not move accuracy, so there is no effect for a mechanism to explain, and we
have not established that the same geometry carries over. The mechanism result
and the decomposition result agree about which benchmark has an operator effect;
they do not license a claim that the mechanism is universal.

\section{Discussion and Conclusion}

Mask-aware attention is not one mechanism but two, and they do not behave alike.
The operator that suppresses unobserved steps costs up to 6.5 macro-F1 points on
one benchmark, does nothing measurable on a second, and gains up to 5.7 points
on a third. The input channels it usually arrives with---a mask indicator,
elapsed time, a calendar encoding---behave in yet another way, and are
themselves not one thing: the calendar and the elapsed-time channel have
opposite signs and largely cancel, which is why the block they form read as a
single weak effect in the conventional comparison.

This matters beyond our three datasets. The comparison this literature
routinely makes, a masked model against an unmasked one, changes both components
at once and therefore cannot attribute what it measures. On two of our three
benchmarks the reported sign belongs to the input channels rather than to the
operator the comparison is named for. The capacity-matched control that
separates them costs one extra training configuration, needs no new data, and is
self-checking: at zero removal it must return exactly zero, and it does.

The three benchmarks disagree about the operator's sign, and we have argued that
the disagreement is the finding rather than a nuisance. It also bounds what can
be claimed. A practitioner choosing whether to mask cannot read a direction off
this literature, because the direction is a property of the benchmark; what can
be read off is that the decision should be evaluated on the target dataset with
the two components separated, and that the separation is cheap.

The sequence-length account that motivated this work does not survive. It came
from stratifying test shards by natural acquisition count, which is
observational; under a within-sample intervention the operator's contribution
shows no trend with the number of surviving observations at all.

A class-separability moderator we reported earlier does not generalise either.
Using two class pairs chosen by clean per-class F1, it showed masking helping a
separable pair and hurting a confusable one. Applied continuously across all 15
TimeSen2Crop classes, the effect on a class does not track its separability
($r = +0.52$ at 30\% removal, $-0.53$ at 90\%); at 70\% every class is harmed,
with no relation to how separable it is. We report the negative result rather
than the two pairs.

\paragraph{Limitations.}
The most serious limitation is geographic. All three benchmarks cover
temperate European agriculture, and two of them are French. We acquired and
discarded two non-European candidates, and the reason was not access but
transferability: no openly downloadable crop-type benchmark we examined supports
a continental-scale spatial holdout, because their published splits are
within-region splits. The third benchmark therefore adds a configuration, not a
regime, and the question of whether the decomposition behaves differently under
tropical or smallholder agriculture remains open. TimeSen2Crop's test set is
additionally balanced by its sampling cap, so absolute macro-F1 should not be
compared across datasets.

Our baselines are confined to one model family per attention type so that
capacity and mechanism stay separable; comparison against state-of-the-art
architectures is a separate question. The second attention family was run on
one benchmark only, so the operator's cross-dataset variation is established for
L-TAE alone. The alignment control covers the two benchmarks where the calendar
does something; on TimeSen2Crop the channel is inert and there is nothing to
control for. Finally, a decomposition is not a causal claim about the loss
landscape: we manipulate the architecture and measure the outcome, and we
report where the mechanism we can measure does and does not close.

\section*{Declarations}
\textbf{Data.} TimeSen2Crop (Zenodo)~\cite{timesen2crop},
BreizhCrops~\cite{breizhcrops} and PASTIS~\cite{utae} are public. Code,
configuration files and the complete run log are released, and every number in
this paper is reproduced by the commands listed in the accompanying results
manifest. \textbf{Ethics.} Public data only; no human subjects.
\textbf{Conflicts.} None. \textbf{Funding.} None. \textbf{AI.} AI tools
assisted coding, data retrieval and language editing; all measurements come
from the released code and log.

\begin{thebibliography}{1}

\bibitem{breizhcrops}
M.~Ru{\ss}wurm, C.~Pelletier, M.~Zollner, S.~Lef\`evre, and M.~K\"orner,
``BreizhCrops: A time series dataset for crop type mapping,'' in
\emph{Int. Arch. Photogramm. Remote Sens. Spatial Inf. Sci.},
vol.~XLIII-B2-2020, 2020, pp.~1545--1551.

\bibitem{timesen2crop}
G.~Weikmann, C.~Paris, and L.~Bruzzone, ``TimeSen2Crop: A million labeled
samples dataset of Sentinel-2 image time series for crop-type classification,''
\emph{IEEE J. Sel. Topics Appl. Earth Observ. Remote Sens.}, vol.~14,
pp.~4699--4708, 2021.

\bibitem{utae}
V.~Sainte Fare Garnot and L.~Landrieu, ``Panoptic segmentation of satellite
image time series with convolutional temporal attention networks,'' in
\emph{Proc. IEEE/CVF Int. Conf. Computer Vision (ICCV)}, 2021,
pp.~4872--4881.

\bibitem{ltae}
V.~Sainte Fare Garnot and L.~Landrieu, ``Lightweight temporal self-attention
for classifying satellite image time series,'' in \emph{Advanced Analytics and
Learning on Temporal Data (AALTD)}, Springer, 2020, pp.~171--181.

\bibitem{vaswani}
A.~Vaswani \emph{et al.}, ``Attention is all you need,'' in \emph{Advances in
Neural Information Processing Systems 30 (NIPS)}, 2017, pp.~5998--6008.

\bibitem{tsvit}
M.~Tarasiou, E.~Chavez, and S.~Zafeiriou, ``ViTs for SITS: Vision transformers
for satellite image time series,'' in \emph{Proc. IEEE/CVF Conf. Computer
Vision and Pattern Recognition (CVPR)}, 2023, pp.~10418--10428.

\bibitem{alise}
I.~Dumeur, S.~Valero, and J.~Inglada, ``Paving the way toward foundation
models for irregular and unaligned satellite image time series,'' \emph{IEEE
Trans. Geosci. Remote Sens.}, vol.~63, pp.~1--21, 2025.

\bibitem{jointfrp}
Y.~Wang \emph{et al.}, ``A joint learning framework with feature
reconstruction and prediction for incomplete satellite image time series in
agricultural semantic segmentation,'' arXiv:2505.19159, 2025.

\bibitem{mae}
Liu \emph{et al.}, ``Crop spectral data reconstruction and classification based
on temporal masked autoencoder,'' in \emph{EGU General Assembly}, EGU26-8646,
2026. (Conference abstract.)

\bibitem{shaping}
G.~Castiglioni, N.~Isla, C.~Buc, J.~Castillo-Navarro, S.~Lef\`evre, and
V.~Barriere, ``Shaping fine-tuning of geospatial foundation models: Effects of
label availability and temporal resolution,'' in \emph{Proc. TerraBytes ICML
Workshop}, PMLR, vol.~292, 2025, pp.~81--96.

\bibitem{eurosat}
P.~Helber, B.~Bischke, A.~Dengel, and D.~Borth, ``EuroSAT: A novel dataset and
deep learning benchmark for land use and land cover classification,'' \emph{IEEE
J. Sel. Topics Appl. Earth Observ. Remote Sens.}, vol.~12, no.~7,
pp.~2217--2226, 2019.

\bibitem{padding}
M.~Dwarampudi and N.~S.~S. Reddy, ``Effects of padding on LSTMs and CNNs,''
arXiv:1903.07288, 2019.

\bibitem{sitsreview}
A.~Zhang, Z.~Zhang, K.~Shi, and P.~Tang, ``Benchmark datasets for satellite
image time series classification: A review,'' \emph{Remote Sensing}, vol.~18,
no.~10, art.~1581, 2026.

\end{thebibliography}

\end{document}
